\chapter{Discussion}
\label{sec:discussion}

\section{The Model, Not the Architecture, Is the Bottleneck}

The central argument of this thesis was that an agent which searches until its
evidence suffices, checks every claim, and abstains otherwise, would turn confident
errors into stated gaps. With the local 7B model, the first half of this happened and
the second did not: the agent was wrong far less often than single-pass, but it held
back easy questions about as readily as hard ones, so it also lost correct answers.

Running the identical graph, prompts and settings with a stronger model separates the
two explanations. With Gemini, the Sufficiency Gate was convinced about five times as
often, the Faithfulness Gate rejected far fewer claims, and --- the decisive
observation --- the questions it held back were the ones single-pass also got wrong.
The gates behave as designed when the model behind them can judge sufficiency and
read a table or chart reliably. A 4-bit 7B model cannot yet do either well enough:
its judgements of sufficiency are not calibrated to difficulty, and half of its drafted
claims fail verification. The architecture is therefore not refuted by the 7B
results; its benefit is gated by model capacity. This also suggests a practical design
for constrained settings: a small local model to read and write, and a stronger model
only for the two judgements.

\section{What the Development Error Analysis Showed}

Before the test run, the full agent's answers on the 20 development questions were
analysed by hand (Table~\ref{tab:dev-errors}). The two most frequent causes were fixed,
after which tuning stopped, as planned.

\begin{table}[htbp]
\centering
\small
\begin{tabular}{p{9.5cm}c}
\toprule
Cause (13 questions not answered correctly, prompt v6) & Count \\
\midrule
G2 too strict: a faithful quote rejected for PDF text quirks (``in- ference'', ``0 . 5\%'') & 4 \\
G1 too lenient: ``sufficient'' with a fact that does not answer the question & 2 (+1) \\
The writer misread the evidence or answered vaguely & 3 \\
The wrong table or figure was chosen & 2 \\
A table with a two-level header was read wrongly & 1 \\
A computed ratio that cannot be checked against the evidence & 1 \\
\bottomrule
\end{tabular}
\caption{Causes of the full agent's errors on the development questions, before the
two fixes (normalised quote matching; a Sufficiency Gate that must name the answering
fact before its verdict).}
\label{tab:dev-errors}
\end{table}

The second fix overshot: the stricter gate produced no wrong answer on development but
13 abstentions, and on test it abstained on more than three questions in four. One
calibration step was tried and brought wrong answers back without reducing the
abstentions, so the stricter gate was kept and the best-effort rule
(Section~\ref{sec:best-effort}) was added instead. The lesson generalises: a gate that
decides both \emph{when to stop searching} and \emph{whether to answer} is asked to
make two different judgements with one threshold.

\section{Measuring Grounding}

IoU against a hand-drawn box is the natural measure of box precision, and it was
fixed before the evaluation. It turned out to measure the agreement of two drawing
conventions as much as the precision of a proof: a parsed table cell spans its
column, a human box hugs the number, and a human asked to box ``the answer'' to a
comparison may box a whole row. The pointing game is insensitive to these conventions
and was therefore added; it is reported beside IoU, not instead of it, and it was added
after the IoU results were seen. A cleaner protocol would fix the unit of annotation in
advance --- for tables, whole cells --- so that both sides draw at the same
granularity.

\section{Practical Constraints}

The system was built and evaluated under three constraints that shaped it. One
consumer GPU with 12~GB of memory holds one large model at a time, so the page index
was built before the language model started and ColQwen2 ran at question time only
from cached vectors. Daily power cuts made every long job resumable and every file
write atomic, and the model-reply cache, kept under version control, made a power cut
cost at most one question. The PC was shared with other experiments, so every model
load first checked free memory and GPU memory, after an early crash of another job
when memory ran out. The stronger-model comparison used a free tier and took three
sittings because of rate limits and busy servers; the last 25 questions used a second
free-tier key, with the model, prompts and settings unchanged.

\section{Threats to Validity}

\begin{itemize}
    \item \textbf{The judge is a model}, and the same model as the evaluated local
    systems. Its agreement with a human grader was substantial (kappa \nJudgeKappa{} on
    25 answers), and it erred towards strictness, but 30 graded answers is a small check.
    \item \textbf{One annotator}, who is a co-author and knew the project's aims,
    produced the boxes and the grades; the first quality-control round was redrawn by
    the first author. The 50 annotated questions bound the precision of every grounding
    number.
    \item \textbf{Post-hoc decisions.} The best-effort rule was introduced after a first
    test run, on development data only, and all systems were then re-run; the pointing
    game was added after the IoU results. Both are disclosed, and the original
    measures are reported.
    \item \textbf{Scope.} 150 questions from 31 computer-science papers; retrieval
    within a single known paper rather than across a corpus; a 4-bit model; one run per
    system at temperature 0. The multi-hop subset has ten questions.
    \item \textbf{MAP} was measured by an automatic stand-in that undercounts correct
    citations of supporting text.
    \item \textbf{Verifiability by humans} (RQ3's second half) was argued from the
    precision of the proofs, not measured with a user study.
\end{itemize}
